% Ad Familiares 9.8 --- headnote
% ad-familiares-09-08, 11 or 12 July 45 BC, Tusculum

\textit{Cicero to M. Terentius Varro, written from the Tusculan villa
on 11 or 12 July 45 BC (Perseus: \textit{v aut iv Id. Quint.}). This
is the dedicatory letter for the \textit{Academica} in its
two-book recension --- the so-called \textit{Academici libri} --- in
which Cicero has cast a dialogue at his Cuman villa with Varro and
Atticus (``Pomponius'') as the interlocutors. Cicero has assigned to
Varro the role of Antiochus of Ascalon, whose Old-Academic position
Varro is known to have favored, and has reserved for himself the
sceptical New-Academic part of Philo of Larissa, his own teacher.
The conceit of the letter is that of a \textit{munus} owed and
returned: Varro has long promised Cicero a literary gift of his own,
and Cicero --- impatient at the delay, but interpreting it
charitably as care rather than neglect --- offers his philosophical
dialogue in advance as the matching present.}

\textit{The ``four reminders not overly modest'' that Cicero says
he has sent are the four books of the \textit{Academica} themselves,
personified as messengers from the ``younger Academy'' --- a wry
play on the brashness of the New Academy he is sending Varro on his
behalf. The closing turn is darker: in better times the two old
scholars could pursue these studies among other honorable
occupations, but with the Republic gone there is nothing else left
to live for. The reference to Varro's ``move and purchase'' --- a
property he is acquiring --- closes on the ordinary note of friendly
business at the end of a programmatic letter of high importance.}
